% ch01-intro.tex – Chapter 1: introduction and background (round 2, Writer).
\ChapterMission{המשימה: להציג את הבעיה, את התובנה שממנה יוצא הפרויקט ואת מבנה הדוח}
\ChapterIcon{pedestrian}
\chapter{מבוא ורקע}
\begin{scenebox}
מוקדם בבוקר, והמדרכה ריקה. גבר אחד עומד לבדו מול איש אדום שנדמה כאילו לא יתחלף לעולם.
הכביש שקט, אין מכונית ואין מי שיביט בו. הוא מסתכל ימינה, שמאלה ושוב ימינה,
ויורד מן המדרכה.
\end{scenebox}

הרגע הזה, שבו הולך רגל יחיד מחליט שההמתנה אינה שווה את מחירה, הוא נקודת המוצא של הדוח.
הפרק מציג את בעיית בטיחות הולכי הרגל בעולם ובישראל, את התובנה של מאמר הדגל, את הפרויקט המוצע ואת מבנה הדוח.

\section{בטיחות הולכי רגל בעולם}
לפי הדוח העולמי של ארגון הבריאות העולמי (\en{WHO}) על בטיחות בדרכים, מספר ההרוגים השנתי בתאונות דרכים ירד מעט, אך עדיין עומד על כ-\num{1.19} מיליון בני אדם בשנה \cite{who2023}.
תאונות דרכים הן סיבת המוות המובילה של ילדים וצעירים בני \range{5}{29}, והולכי רגל, רוכבי אופניים ורוכבי אופנועים מהווים יחד יותר ממחצית ההרוגים, רובם במדינות בעלות הכנסה נמוכה ובינונית \cite{who2023}.
היעד העולמי הוא להפחית לפחות במחצית את מספר ההרוגים והפצועים עד שנת \num{2030} \cite{who2023}.
הולכי הרגל הם אפוא קבוצה מרכזית בכל מאמץ לעמוד ביעד זה, משום שאין להם כל מיגון מול כלי הרכב.

\section{התמונה בישראל}
בישראל, מדינה בעלת הכנסה גבוהה, משקלם של הולכי הרגל בין הרוגי הדרכים ניכר.
לפי פרופיל המדינה של ארגון הבריאות העולמי, בשנת \num{2021} דווחו בישראל \num{364} הרוגים בתאונות דרכים, ואומדן הארגון הוא \num{4.2} הרוגים לכל \num{100,000} תושבים \cite{who2023israel}.
הולכי הרגל היוו כ-\pct{27} מן ההרוגים המדווחים, כלומר כמעט אחד מכל ארבעה \cite{who2023israel}.
במונחים מוחלטים מדובר בקרוב למאה הולכי רגל בשנה. זהו חישוב של המחברים, \pct{27} מתוך \num{364}, ולא נתון שפורסם במקור.
% LEADER: the derived value (0.27 × 364 ≈ 98) is not in data/recomputed.json yet. Once registered, show it as \num{98} here.
עוד עולה מן הפרופיל ש-\pct{77} מן ההרוגים המדווחים הם גברים \cite{who2023israel}. נתון זה מתיישב עם הפער המגדרי שנמצא בחצייה באדום, אך הוא נתון רקע בלבד ואינו ראיה לגבי חצייה באדום.

\begin{statrow}[3]
  \stattile[SignalRed]{\num{1.19}}{מיליון הרוגים בתאונות דרכים בעולם בשנה}
  \stattile[SignalAmber]{\pct{27}}{מן ההרוגים בישראל בשנת \num{2021} היו הולכי רגל}
  \stattile[SignalGreen]{\pct{13.5}}{מן המגיעים באדום חצו באדום במעבר שנצפה בתל אביב}
\end{statrow}

\section{התובנה ממאמר הדגל}
מאמר הדגל של הדוח בחן \num{1,392} הולכי רגל במעבר חצייה מרומזר בתל אביב, ומצא ש-\pct{13.5} ממי שהגיעו במופע האדום חצו באדום \cite{rosenbloom2009}.
התובנה המרכזית אינה נוגעת למספר המפרים אלא לנסיבות: הולך רגל שעמד לבדו על המדרכה חצה באדום יותר מהולך שעמד עם אחרים, וכל אדם נוסף שהמתין על המדרכה הקטין את הסיכויים לחצייה באדום \cite{rosenbloom2009}.
גברים חצו באדום יותר מנשים.
המחבר מסביר את הממצא בשליטה חברתית: הקבוצה הממתינה מרסנת את חבריה. הוא ממליץ לנצל זאת בקמפיינים, ולקצר את משך האדום בשעות מסוימות של היום \cite{rosenbloom2009}.
המאמר גם מציין שזמן ההמתנה לא נמדד ושהתצפית נערכה באתר אחד בלבד. פרק \num{4} מנתח את המאמר ואת נתוניו באופן ביקורתי.

\section{הפרויקט בפסקה אחת}
הפרויקט המוצע הוא מעבר חצייה חכם, שפועל על הגורם שמאמר הדגל זיהה: הולך הרגל הממתין, ובמיוחד הממתין היחיד.
חישה בשפת המדרכה מזהה הולכי רגל ממתינים, את מספרם ואת זמן המתנתם, ונותנת עדיפות להולך יחיד; הבקר מקצר בהתאם את מופע האדום להולכי הרגל, בגבולות בטוחים שנקבעים בפרק הדרישות. זהו האדום המסתגל (\en{adaptive red}).
תצוגת ספירה לאחור (\en{countdown}) פועלת רק במופע הפינוי (\en{clearance interval}), כלומר בזמן שהאות להולכי הרגל מהבהב, כפי שהתקן האמריקני מחייב \cite{fhwa2025mutcd}; בזמן האדום אין ספירה.
לצד המערכת פועל קמפיין שליטה חברתית, באתר ובתקשורת, המציג את ההמתנה על המדרכה כנורמה.

\subsection{שינוי כיוון מבוסס ראיות}
הרעיון המקורי של הפרויקט היה תצוגת ספירה לאחור שמראה לממתין כמה זמן נותר לאדום.
סקירת המקורות הראתה שני דברים. ראשית, התקן מתיר ספירה להולכי רגל רק במופע הפינוי, ולא בזמן האדום (סעיף \en{4I.04}) \cite{fhwa2025mutcd}. שנית, במחקר שהשווה צמתים ליד בתי ספר, ספירה בזמן האדום הלכה יחד עם יותר הפרות ויותר יציאות מוקדמות בקרב ילדים \cite{fu2016}.
בעקבות זאת הועברה הספירה למופע הפינוי, והחישה הופנתה לקיצור האדום עצמו: במקום להראות לממתין כמה זמן נותר, המערכת מקצרת את ההמתנה.
שינוי הכיוון הוא חוזקה של התכנון ולא חולשה: ההחלטה נגזרה מן הראיות ומן התקן, ולא מן ההעדפה הראשונית.
פרק \num{6} מציג את הספירה בזמן האדום כחלופה הראשונה ומראה כיצד היא נדחית במטריצת ההחלטה המשוקללת.

\section{מבנה הדוח}
הדוח בנוי מעשרה פרקים, ואחרי המבוא באים תשעה:
\begin{itemize}
\item פרק \num{2}, הגדרת הבעיה, שאלת המחקר והיעדים: הבעיה במספרים, ההשערות ויעדי התכן.
\item פרק \num{3}, מפרט הדרישות: דרישות פונקציונליות ולא פונקציונליות, כולל הגבולות הבטוחים לקיצור האדום, כל אחת עם יעד ושיטת אימות.
\item פרק \num{4}, סקירת ספרות ומאמר הדגל: סקירה תמטית, מטריצת ספרות, ניתוח ביקורתי של נתוני מאמר הדגל ופער הידע.
\item פרק \num{5}, חקר שוק ומוצרים קיימים: השוואה טכנית של תצוגות ספירה, גלאי הולכי רגל ומעברים חכמים.
\item פרק \num{6}, חלופות תכן ומטריצת החלטה: שלוש חלופות, משקלות מנומקים ובחירה.
\item פרק \num{7}, ארכיטקטורת המערכת וחדשנות: מן החישה, דרך הבקר, אל התצוגה ושכבת הקמפיין והנתונים, כולל פרטיות כבר בתכנון.
\item פרק \num{8}, תוכנית עבודה: חבילות עבודה, מטריצת אחריות ותרשים גאנט.
\item פרק \num{9}, ניתוח סיכונים ותוכנית הפחתה.
\item פרק \num{10}, סיכום: הממצאים, מוכנות הפרויקט והצעדים הבאים.
\end{itemize}
